\chapter*{Conclusion generale}
\addcontentsline{toc}{chapter}{Conclusion generale}

Ce projet avait pour objectif la conception et le developpement de \textbf{FinPulse},
une plateforme intelligente d'analyse financiere basee sur les donnees reglementaires
de la SEC et sur les techniques d'intelligence artificielle. A travers ce travail,
nous avons cherche a repondre a une problematique centrale : comment aider
l'investisseur a mieux evaluer la coherence entre une idee d'investissement, le
discours public d'une entreprise et les risques officiellement declares dans ses
rapports financiers.

La solution proposee s'appuie sur l'exploitation des rapports 10-K, qui constituent
une source fiable et officielle pour l'analyse des entreprises cotees. Ces documents,
bien que riches en informations, sont souvent difficiles a analyser manuellement en
raison de leur volume, de leur complexite et de leur langage technique. FinPulse
apporte une reponse a cette difficulte en automatisant l'ingestion, l'extraction, la
vectorisation et l'analyse semantique des donnees financieres.

Sur le plan fonctionnel, la plateforme integre plusieurs modules complementaires.
L'assistant d'analyse permet a l'utilisateur de formuler des questions ou des idees
d'investissement en langage naturel et d'obtenir des reponses argumentees a partir des
documents officiels. Le moteur de strategie permet de produire une analyse plus
structuree, en mettant en evidence les preuves, les contradictions et les signaux de
risque. La surveillance automatisee permet de suivre les strategies sauvegardees et de
declencher des alertes lorsque des changements significatifs sont detectes. Enfin, la
watchlist en temps reel offre une vision synthetique de l'evolution des entreprises
suivies, a travers des indicateurs tels que le score NCI, les variations boursieres,
les actualites pertinentes et le sentiment shift.

Sur le plan technique, le projet repose sur une architecture moderne, modulaire et
evolutive. L'utilisation d'Angular pour l'interface utilisateur, de Spring Boot pour
l'orchestration metier, de Python/FastAPI pour le traitement des donnees et de
PostgreSQL avec PGVector pour le stockage structure et vectoriel permet de separer
clairement les responsabilites entre les differents composants. L'integration de Redis
contribue egalement a ameliorer les performances par la mise en cache des donnees
frequemment consultees.

Ce travail nous a permis de mettre en pratique plusieurs competences acquises durant
notre formation, notamment la conception UML, l'architecture logicielle, le
developpement web, la gestion des bases de donnees, l'integration d'API, ainsi que
l'utilisation de techniques d'intelligence artificielle appliquees a l'analyse de
documents. Il a egalement mis en evidence l'importance de la qualite des donnees, de
la tracabilite des resultats et de la securite dans le developpement d'une plateforme
d'aide a la decision.

Malgre les apports de la solution, certaines limites peuvent etre relevees. Le systeme
depend de la qualite et de la disponibilite des sources externes, notamment les
documents SEC, les donnees de marche et les actualites financieres. De plus,
l'interpretation automatique des documents financiers reste un domaine complexe, ou la
precision des resultats depend fortement des modeles utilises, des prompts definis et
de la qualite des donnees d'apprentissage ou de reference. Une validation humaine peut
donc rester necessaire pour les decisions financieres sensibles.

Plusieurs perspectives d'amelioration peuvent etre envisagees. Il serait possible
d'etendre la plateforme a d'autres types de documents reglementaires, comme les
rapports 10-Q ou les formulaires 8-K, afin d'enrichir l'analyse. L'ajout de nouveaux
modeles de prediction, de tableaux de bord plus avances, d'indicateurs sectoriels et
d'un systeme de recommandation personnalise permettrait egalement d'ameliorer la
valeur analytique de la solution. Enfin, une phase de tests a plus grande echelle
pourrait permettre d'evaluer la robustesse, la performance et la pertinence des
analyses produites.

En conclusion, FinPulse constitue une base solide pour une plateforme d'aide a la
decision financiere intelligente, explicable et evolutive. En combinant les donnees
officielles, les architectures modernes et l'intelligence artificielle, le projet
montre comment les technologies actuelles peuvent contribuer a rendre l'analyse
financiere plus accessible, plus transparente et mieux adaptee aux besoins des
investisseurs.
